\chapter{L'orario come oggetto strutturale}
\label{ch:orario_oggetto_strutturale}

\epigraph{``First, solve the problem. Then, write the
code.''}{John Johnson}

\lettrine[lines=3]{P}{rima} di descrivere come $\pi$Tantum
\emph{costruisce} un orario, vale la pena descrivere
\emph{cosa \`e} un orario --- visto come oggetto strutturato.
Le entit\`a coinvolte, le relazioni che le legano, i
vincoli ricorrenti che intervengono. Questa lente
``strutturale'' \`e quella che ha guidato la progettazione
del modello dati di $\pi$Tantum (cap.~\ref{ch:modello_dati}).

\section{Le entit\`a}

Una scuola italiana, ai fini del timetabling, si descrive
con sette entit\`a principali.

\subsection{Docente}

Il \emph{docente} ha un'anagrafica (cognome, nome, eventuale
nickname), una o pi\`u \emph{classi di concorso} (es. A027
matematica e fisica, A019 storia e filosofia), un \emph{monte
ore massimo} settimanale (tipicamente 18 per il ruolo
ordinario), un \emph{giorno libero contrattuale}, e una
\emph{matrice di disponibilit\`a} 6$\times$6 in cui ogni cella
ha uno dei cinque stati (allowed, soft, hard, preferred,
enforced).

Pu\`o avere preferenze aggiuntive: una \emph{matrice di
preferenza per le aule} (in quali aule vorrebbe insegnare,
con i 5 stati), un \emph{punteggio di graduatoria} per il
criterio Anzianit\`a, un elenco di \emph{tre giorni preferiti
per essere libero}.

\subsection{Classe}

La \emph{classe} ha un nome (es. ``3A scientifico''), un
\emph{anno di corso} (1--5), un \emph{indirizzo} (rimanda
all'entit\`a Indirizzo), un \emph{numero di studenti}, una
\emph{home class} (l'aula di norma), e una propria matrice
di disponibilit\`a 6$\times$6 (es. la 3A non fa lezione il
sabato perch\'e in gita una volta al mese).

Pu\`o avere un \emph{max\_hours\_per\_day} (limite di ore
giornaliere, tipicamente 6 ma talora 5 o 7 per casi
speciali), un \emph{peso SOFT sulla 6\textsuperscript{a} ora}
(quanto si vuole evitarla), e tre \emph{giorni preferiti
liberi}.

\subsection{Materia}

La \emph{materia} ha un nome (Matematica, Italiano, Storia,
Religione, ...), una \emph{matrice di compatibilit\`a con le
aule} (es. fisica deve stare in lab\_fisica, religione pu\`o
stare ovunque), un peso che indica quanto la materia \`e
``concettualmente pesante'' (per il vincolo di ``no due ore
pesanti consecutive''), e una \emph{matrice subject-group-weight}
che dice quali classi di concorso possono insegnarla.

\subsection{Aula}

L'\emph{aula} ha un nome (es. ``Aula 12'' o ``Lab Fisica 1''),
un \emph{tipo} (\texttt{aula\_normale}, \texttt{lab\_fisica},
\texttt{lab\_chimica}, \texttt{lab\_informatica},
\texttt{palestra}, \texttt{aula\_magna}, ecc.), una
\emph{capienza}, un flag \emph{multi\_class} (se pu\`o
ospitare pi\`u classi contemporaneamente, es. l'aula magna),
una \emph{matrice di disponibilit\`a} 6$\times$6 (per
indisponibilit\`a tipo manutenzione), una \emph{matrice di
preferenze classe-by-classe}, e una lista di \emph{tag}
(stringhe libere come \texttt{matematica}, \texttt{piano\_terra},
\texttt{accessibile\_disabili}).

\subsection{Indirizzo}

L'\emph{indirizzo} (curriculum) corrisponde all'``ordinamento''
italiano (Liceo Classico, Scientifico, Linguistico, ITIS,
IPSIA, IPS, ITER, ITT, ...). Per ogni anno e per ogni
materia, l'indirizzo specifica il \emph{numero di ore
settimanali}. Un'eventuale lista di \emph{vincoli logici
indirizzo-specifici} (es. ``in 1\textsuperscript{o} anno
scientifico la matematica non sta dopo la 4\textsuperscript{a}
ora'').

\subsection{Studente}

Lo \emph{studente} ha un'anagrafica, appartiene a una classe
di base, pu\`o appartenere a uno o pi\`u \emph{gruppi
articolati} (es. seconda lingua, IRC vs alternativa,
recupero), ha eventuali \emph{tag} (\texttt{BES},
\texttt{DSA}, \texttt{debito\_matematica\_4}). I tag sono
metadato passivo per filtri e raggruppamenti, non sostituiscono
i gruppi.

\subsection{Gruppo articolato}

Il \emph{gruppo articolato} \`e un insieme di studenti che si
incontrano in slot comuni per un'attivit\`a specifica (es.
``IRC della 1A+1B'', ``Recupero matematica del biennio'',
``Spagnolo seconda lingua del triennio''). I membri possono
provenire da qualunque combinazione di classi (semantica
``Tipo C''). Un gruppo richiede slot dedicati e un'aula
specifica.

\section{Le relazioni}

Le entit\`a sopra sono collegate da una rete di relazioni che
sono il ``substrato'' del timetabling.

\begin{itemize}
\item \textbf{Cattedra} (Docente $\to$ Classe $\to$ Materia
  $\to$ ore): l'unit\`a fondamentale. Indica che il docente $D$
  insegna $h$ ore della materia $M$ alla classe $C$ in
  un anno.
\item \textbf{Compresenza}: due (o pi\`u) docenti che devono
  essere nello stesso slot stessa classe stessa materia.
  Tipicamente teoria + laboratorio nei tecnici.
\item \textbf{Lezione}: un'istanza concreta in un certo
  slot, con un'aula assegnata. \`E ci\`o che la pipeline
  produce.
\item \textbf{Indisponibilit\`a (HARD)}: la matrice 6$\times$6
  con celle rosse del docente/classe/aula.
\item \textbf{Vincolo logico DSL}: un'espressione del DSL
  generico (cap.~\ref{ch:metodo_dsl}) che lega entit\`a in
  qualunque combinazione.
\end{itemize}

\section{I vincoli ricorrenti}

Alcuni vincoli ritornano in ogni scuola italiana, anche se
con sfumature diverse.

\begin{itemize}
\item \textbf{Copertura}: ogni materia deve ricevere il suo
  monte ore in ogni classe. HARD.
\item \textbf{Non sovrapposizione}: docente, classe, aula
  in al pi\`u una lezione per slot. HARD.
\item \textbf{Compatibilit\`a classe di concorso}: solo
  docenti della classe giusta possono insegnare una
  materia. HARD.
\item \textbf{Compatibilit\`a aula--materia}: fisica solo in
  lab\_fisica, ginnastica solo in palestra, etc. HARD o
  PREFERRED a seconda della scuola.
\item \textbf{Giorno libero contrattuale}: un giorno di
  riposo settimanale per ogni docente. HARD per CCNL.
\item \textbf{Compresenze}: due lezioni esattamente
  nello stesso slot stessa classe. HARD.
\item \textbf{No buchi nelle classi}: nessuna ora vuota nel
  mezzo della giornata della classe. HARD nelle scuole
  pi\`u rigide, SOFT in quelle pi\`u flessibili.
\item \textbf{Sesta ora}: si vorrebbe evitare la sesta ora,
  o renderla solo per certe materie. SOFT.
\item \textbf{No buchi docenti}: simile, ma per i docenti.
  SOFT.
\item \textbf{Distribuzione settimanale}: una stessa
  materia non concentrata in due giorni consecutivi. SOFT.
\item \textbf{Coppie di ore consecutive}: certe materie
  (es. ginnastica, lab) richiedono blocchi di 2 ore.
  PREFERRED o HARD.
\end{itemize}

\section{Il caso italiano: peculiarit\`a}

\begin{casostudiobox}
\textbf{IIS ``Galileo Ferraris''}, Torino, 28 classi
distribuite su tre indirizzi (ITIS Informatica, ITIS
Elettronica, IPS Servizi Commerciali). Le tre articolazioni
hanno monte ore diversi anche per materie ``comuni'' come
matematica e italiano. I docenti di matematica (A026 e
A027) coprono tutte e tre le articolazioni; i docenti
specifici (es. Sistemi e Reti per Informatica) sono
limitati a un solo indirizzo. Compresenze numerose
(matematica + lab\_inf in Informatica, fisica + lab\_fisica
in Elettronica). Giorni di stage in azienda concentrati
nelle ultime settimane: $\pi$Tantum gestisce attraverso un
\emph{vincolo logico DSL} che svuota i pomeriggi di marted\`i
e gioved\`i delle classi quinte nel periodo aprile--maggio.
\end{casostudiobox}

Il caso italiano si distingue dai benchmark accademici
internazionali (ITC, XHSTT~\parencite{mccollum2010}) per:
\begin{itemize}
\item Indirizzi (curricula) come entit\`a di prim'ordine,
  con monte ore variabile.
\item Classi di concorso esplicite e gerarchiche.
\item Compresenze come vincolo HARD ricorrente.
\item Gruppi articolati cross-class.
\item Vincoli sindacali (giorno libero CCNL, max ore).
\end{itemize}

Tutte queste peculiarit\`a sono codificate nel modello dati
di $\pi$Tantum (cap.~\ref{ch:modello_dati}); $\pi$Tantum
\`e progettato \emph{a partire} da queste, non come
generalizzazione di un sistema accademico.
